
\chapter{Results}
\begin{spacing}{1.3}

\section{Model Performance}

Three architectures were trained, evaluated, and compared on the held-out test set of 1,527 images with patient-aware splitting ensuring no lesion-level leakage. Table~\ref{tab:final} presents the complete evaluation using the clinically-calibrated decision threshold of 0.15.

\begin{table}[H]
\centering
\caption{Final model comparison on held-out test set (1,527 images).}
\label{tab:final}
\begin{tabular}{lrrr}
\toprule
\textbf{Metric} & \textbf{EfficientNet-B0} & \textbf{ResNet50} & \textbf{MobileNetV3-S} \\
\midrule
Parameters          & \textbf{4.3M} & 24M  & 1.5M \\
ROC-AUC             & \textbf{0.91} & 0.84 & 0.84 \\
Sensitivity @ 0.15  & \textbf{0.806} & 0.71 & 0.69 \\
Specificity @ 0.15  & 0.828 & 0.80 & \textbf{0.81} \\
Precision (PPV)     & 0.35 & 0.29 & 0.31 \\
F1 Score            & \textbf{0.49} & 0.41 & 0.43 \\
Accuracy            & 0.826 & 0.79 & \textbf{0.79} \\
Balanced Accuracy   & \textbf{0.817} & 0.76 & 0.75 \\
Inference time (ms) & 6.7 & 14.2 & \textbf{4.1} \\
ONNX Size (MB)      & 16.55 & 94.40 & \textbf{6.13} \\
\bottomrule
\end{tabular}
\end{table}

EfficientNet-B0 achieves the highest ROC-AUC (0.91) on the held-out test set, a 7-point advantage over both ResNet50 and MobileNetV3-Small. At the clinically-calibrated decision threshold of 0.15 (see Section~\ref{sec:threshold}), sensitivity reaches 80.6\% -- meeting the screening target of detecting 4 out of 5 melanomas -- with specificity of 82.8\%. The 4.3M-parameter model is 5.6$\times$ smaller than ResNet50 while delivering superior performance, validating the compound scaling hypothesis that jointly optimising depth, width, and resolution produces simultaneously more accurate and parameter-efficient architectures.

\textbf{Threshold calibration.} The default 0.5 decision threshold yielded sensitivity of only 41.9\% -- missing more than half of melanomas and rendering the system clinically unsafe for screening. Systematic threshold analysis revealed that lowering the threshold to 0.15 recovers sensitivity to 80.6\% while maintaining specificity above 80\%. This trade-off is appropriate for a screening tool: false positives (approximately 17\% of benign cases flagged for review) are manageable, whereas false negatives (missed melanomas) carry unacceptable clinical risk. The system compensates for the increased false-positive rate through the clinical review flag and uncertainty gating mechanisms.

The precision of 0.35 reflects the severe 8:1 class imbalance in HAM10000. Even a well-calibrated model generates many false positives when the base rate of melanoma is only 11\% -- a known phenomenon in medical screening. This is precisely why the system includes clinical review flags rather than attempting autonomous clinical decisions. ResNet50 achieved AUC 0.84 with primary weakness in sensitivity (0.71 versus 0.806), missing approximately 29\% of melanomas. MobileNetV3-Small (AUC 0.84, 1.5M parameters) performs comparably to ResNet50 despite being 16$\times$ smaller, providing a viable fallback for ultra-low-resource scenarios.

\section{Optimization Results}

\label{sec:optimization}
INT8 quantisation was investigated across five strategies: dynamic QInt8, dynamic QUInt8, static QDQ with entropy calibration, static QDQ with min-max calibration, and per-channel QDQ. \textbf{All five strategies produced unacceptable accuracy degradation} on EfficientNet-B0, with AUC drops ranging from 10 to 15 percentage points (0.91 down to 0.76--0.81). This is consistent with documented limitations of quantising architectures that use SiLU/Swish activation functions and depthwise separable convolutions -- both core features of EfficientNet.

\begin{table}[H]
\centering
\caption{INT8 quantisation results: all strategies fail for EfficientNet-B0.}
\label{tab:quant_results}
\begin{tabular}{lccc}
\toprule
\textbf{Strategy} & \textbf{Size (MB)} & \textbf{AUC} & \textbf{$\Delta$ AUC} \\
\midrule
FP32 (baseline)       & 16.55 & 0.91  & -- \\
Dynamic QInt8         & 4.97  & 0.80  & $-$0.11 \\
Dynamic QUInt8        & 4.97  & 0.76  & $-$0.15 \\
Static QDQ (entropy)  & 4.48  & 0.81  & $-$0.10 \\
Static QDQ (min-max)  & 4.48  & 0.78  & $-$0.13 \\
Per-channel QDQ       & 4.48  & 0.80  & $-$0.11 \\
\bottomrule
\end{tabular}
\end{table}

\textbf{Recommendation:} Deploy FP32 ONNX (16.6~MB). This size is practical for all target edge devices (Raspberry Pi~4 has 2--8~GB RAM, Jetson Nano has 4~GB) and preserves the full diagnostic accuracy. The 10--15\% AUC penalty of INT8 is clinically unacceptable for a screening tool where sensitivity is paramount. This finding is preserved as a research contribution: EfficientNet-B0, despite its parameter efficiency, is fundamentally incompatible with post-training INT8 quantisation at medically-required accuracy levels. Future work could explore quantisation-aware training (QAT) from scratch or alternative architectures with quantisation-friendly activation functions (e.g., MobileNetV3's hard-swish).

\section{Data Validation Results}

All 7 Great Expectations expectations passed against the HAM10000 dataset. The 100\% image integrity check -- parallelised OpenCV loading across 8 threads -- confirmed zero corrupted, truncated, or missing JPEG files among 10,015 images. Zero patient leakage across train/val/test splits was confirmed through set-intersection assertions on lesion identifiers, satisfying NFR6. One near-duplicate cross-split collision was identified (perceptual hash Hamming distance = 3, below the threshold of 5) and resolved by removing the duplicate from the validation set.

\begin{table}[H]
\centering
\caption{Great Expectations validation summary.}
\label{tab:ge}
\begin{tabular}{lll}
\toprule
\textbf{Expectation} & \textbf{Result} & \textbf{Observed Value} \\
\midrule
Column existence              & Passed & 7/7 columns present \\
Column types                  & Passed & Correct types \\
Diagnosis in valid set        & Passed & 7 valid codes \\
No nulls in key columns       & Passed & 0 nulls \\
Image identifier unique       & Passed & 10,015 unique \\
Row count = 10,015            & Passed & 10,015 rows \\
Image integrity               & Passed & 10,015/10,015 readable \\
\bottomrule
\end{tabular}
\end{table}

\section{Explainability Analysis}

Grad-CAM saliency maps confirm clinically coherent model attention: heatmaps consistently focus on pigmented lesion boundaries, irregular texture patterns, and colour variegation -- the same features dermatologists assess using the ABCDE criteria (Asymmetry, Border irregularity, Colour variation, Diameter, Evolution). No systematic attention to image corners, hair artefacts, or ruler markings was observed, providing evidence that the preprocessing pipeline (DullRazor hair removal, Macenko colour normalisation) effectively removes spurious correlations.

MC-Dropout uncertainty distributions show clear separability between correct and incorrect predictions. Incorrect predictions exhibit 2--3$\times$ higher epistemic uncertainty and predictive entropy than correct predictions, validating the safety architecture's use of standard deviation exceeding 0.15 as the uncertainty gate threshold. Temperature scaling with optimal $T = 1.13$ reduced Expected Calibration Error (ECE) from 0.072 to 0.034 -- below the 0.10 target -- meaning that when the model outputs a 70\% melanoma probability, the observed melanoma frequency is approximately 70\%.

\section{Test Suite}

The test suite comprises 41 tests all passing on every CI run: 34 unit tests spanning data layer, preprocessing, model, ONNX, schema, and utility components; 7 integration tests covering all 6 API endpoints, prediction caching, and Prometheus metric emission. Code coverage exceeds the 75\% gate. Lint status: 0 Ruff errors, 0 Ruff formatting violations, enforced through pre-commit hooks with CI double-check.

\section{Gradio Frontend Interface}

\label{sec:gradio}
The interactive web application provides four tabs for exploring model capabilities, deployed on Hugging Face Spaces (free tier). The interface was designed with a professional soft teal colour palette using Gradio's Soft theme.

\subsection{Predict Tab}

The Predict tab accepts dermoscopic image upload and displays the melanoma risk probability alongside a Plotly risk gauge with colour-coded zones: green for low risk (0--30\%), yellow for uncertain (30--70\%), red for high risk (70--100\%). Key metrics -- melanoma probability, confidence level, and uncertainty -- are presented below the image-gauge row. A clinical review warning appears when predictions fall near the decision boundary or exhibit high uncertainty.

\begin{figure}[H]
\centering
\includegraphics[width=0.92\textwidth]{../figures/screenshot_predict.png}
\caption{Gradio interface: Predict tab with risk gauge and clinical review flag.}
\label{fig:predict_tab}
\end{figure}

\subsection{Explain Tab}

The Explain tab generates a Grad-CAM heatmap overlay on the uploaded image, highlighting regions that most influenced the model's prediction. Red and yellow zones indicate high-activation areas. The tab reports the prediction result, melanoma probability, and confidence level alongside the heatmap.

\begin{figure}[H]
\centering
\includegraphics[width=0.92\textwidth]{../figures/screenshot_explain.png}
\caption{Gradio interface: Explain tab with Grad-CAM activation heatmap.}
\label{fig:explain_tab}
\end{figure}

\subsection{Model Details Tab}

The Model Details tab presents a comprehensive technical reference: performance metrics (ROC-AUC 0.91, sensitivity 80.6\%, specificity 83.7\%), a three-model comparison table (EfficientNet-B0 vs ResNet50 vs MobileNetV3-Small), the INT8 quantisation limitation documentation, and the five-step preprocessing pipeline description.

\begin{figure}[H]
\centering
\includegraphics[width=0.92\textwidth]{../figures/screenshot_model_details.png}
\caption{Gradio interface: Model Details tab with performance metrics and technical reference.}
\label{fig:model_details_tab}
\end{figure}

\subsection{About Tab}

The About tab provides project context: background on the ENSIAS final-year engineering project, the technical architecture flow (EfficientNet-B0 $\rightarrow$ ONNX Runtime $\rightarrow$ FastAPI $\rightarrow$ Gradio), the full MLOps toolchain, known limitations (HAM10000 Fitzpatrick scale bias, research-only status), and a link to the API documentation.

\begin{figure}[H]
\centering
\includegraphics[width=0.92\textwidth]{../figures/screenshot_about.png}
\caption{Gradio interface: About tab with project context and limitations.}
\label{fig:about_tab}
\end{figure}

\end{spacing}
